\section{Primary Research Skills}
\label{research:PrimaryResearchSkills}
\subsection*{Survey Skills:}

\textbf{Source:} \cite{WritingSurveys}

\begin{itemize}
	\item Always keep in mind the purpose of the survey
	\item Keep questions simple and to the point
	\item Don't leave questions vague to maximise response usefulness\
	\item Don't user leading questions, ensure neutrality
	\item Do not collect unnecessary data (but it can be positive to give users choice for future interaction)
\end{itemize}

\textbf{Source:} \cite{HowWriteGood}
\begin{itemize}
	\item Dont only ask closed ended questions
	
	\item Ensure answer choices are balanced 
	E.G Not:
	1. Extremely helpful
	2. Very helpful
	3. Helpful
	
	But instead:
	
	1. Very helpful
	2. Helpful
	3. Neither helpful nor unhelpful
	4. Unhelpful
	5. Very unhelpful

	\item Don't ask for multiple answers at once (unless ranking items)
	\item Let most questions be optional to answer: a user who is confused / unsure of a mandatory question may abandon the survey when their other insight is still highly useful or (potentially worse) select a random answer polluting the results
\end{itemize}
\subsection*{Interview Skills:}
\textbf{Source:} \cite{HowInterviewSomeone}
\begin{itemize}
	\item whilst this link focuses on interviewing candidates for a job role I still found many of these techniques valuable

	\item Develop a structured list of questions tailored to the role, balancing behavioural (e.g., “Describe a time when…”) and situational (e.g., “What would you do if…”).
	\item  Determine what you need to learn from the interview

	\item  Begin with a friendly introduction and outline the interview structure to manage expectations and reduce nervousness.

	\item Avoid leading or overly complex questions that might confuse the interviewee. Instead, phrase questions clearly and concisely to ensure accurate responses.
	\item Example of a balanced scale:
	1.	Very helpful
	2.	Helpful
	3.	Neither helpful nor unhelpful
	4.	Unhelpful
	5.	Very unhelpful

	\item Use follow-up questions or prompts like “Can you elaborate?” to gain more detailed responses, ensuring a thorough understanding of the skills and experiences.

	\item Practice active listening and take clear, structured notes during the interview. This helps you focus on the responses' without distractions.
\end{itemize}
